\subsection{Physical changes on the site images}
\label{sec:peps_range_parent_transport}

Fix a finite simple graph with ordered vertices $V$, site maps
$A_v:\mathcal V_v\to\mathbb C^d$, and linear maps
$F_v:\mathbb C^d\to\mathbb C^e$. Write $B_v=F_vA_v$,
$\mathcal G_S(A)=\operatorname{im}\mathcal C_S(A)$ for the full
open-region range, and $\mathcal P_{\mathcal R}(A)$ for the common
parent space of a family $\mathcal R=(R_i)_{i\in\mathcal I}$.
Set $F_S=\bigotimes_{v\in S}F_v$ and
$F_{\mathrm{all}}=\bigotimes_{v\in V}F_v$.
The following are consequences of the site inverse in
\cite[Definition~5.1]{Schuch2010PEPS} and the regional parent construction.
The virtual spaces and representations remain fixed.

\begin{theorem}[Inverses on the site images]%
    \label{thm:peps_range_parent_site_inverse}
    \lean{TNLean.PEPS.exists_siteRange_leftInverse,
        TNLean.PEPS.physicalDeform_physicalDeform_eq_of_leftInverse,
        TNLean.PEPS.physicalDeform_rightInverse_on_site_ranges}
    \leanok
    If every $F_v$ is injective on $\operatorname{im}A_v$, there are
    linear maps $L_v:\mathbb C^e\to\mathbb C^d$ such that
    $L_vF_vx=x$ for $x\in\operatorname{im}A_v$.
    For any such choice, $L_vB_v=A_v$, and $F_vL_vy=y$ for
    $y\in\operatorname{im}B_v$.
\end{theorem}
\begin{proof}\leanok
    The injection $F_v:\operatorname{im}A_v\to\mathbb C^e$ has a
    linear left inverse. Compose it with the inclusion of
    $\operatorname{im}A_v$ into $\mathbb C^d$ to obtain $L_v$.
    Applying $L_vF_v$ to each column of $A_v$ gives $L_vB_v=A_v$.
    If $y=B_vz=F_vA_vz$, then $F_vL_vy=F_vA_vz=y$.
\end{proof}

\begin{theorem}[Reconstruction on every open-region range]%
    \label{thm:peps_range_parent_regional_inverse}
    \lean{TNLean.PEPS.regionPhysicalMap_leftInverse_on_regionGroundSpace,
        TNLean.PEPS.regionPhysicalMap_injOn_regionGroundSpace}
    \leanok
    Suppose $L_vF_v$ fixes $\operatorname{im}A_v$ at every vertex.
    For any region $S$ and $\psi\in\mathcal G_S(A)$,
    $L_SF_S\psi=\psi$.
    Consequently site-image injectivity of the $F_v$ implies
    injectivity of $F_S$ on $\mathcal G_S(A)$ for every $S$.
    All virtual boundary conditions are allowed.
\end{theorem}
\begin{proof}\leanok
    \uses{thm:peps_range_parent_site_inverse}
    Write $\psi=\mathcal C_S(A)\theta$ for arbitrary boundary data
    $\theta$. Physical deformation commutes with contraction, so
    \begin{align}
        L_SF_S\mathcal C_S(A)\theta
        =L_S\mathcal C_S(B)\theta
        =\mathcal C_S(LB)\theta
        =\mathcal C_S(A)\theta.
    \end{align}
    Apply $L_S$ to an equality $F_S\psi=F_S\phi$ to obtain
    injectivity. When only site-image injectivity is given, choose
    $L_v$ by Theorem~\ref{thm:peps_range_parent_site_inverse}.
\end{proof}

\begin{theorem}[Product maps fixing all site slices]%
    \label{thm:peps_range_parent_slice_reconstruction}
    \lean{TNLean.PEPS.globalPhysicalMap_eq_self_of_fixes_site_ranges,
        TNLean.PEPS.globalPhysicalMap_leftInverse_of_mem_parent}
    \leanok
    \uses{def:peps_vertex_image_ground_space}
    Let $P_v:\mathbb C^d\to\mathbb C^d$ fix
    $\operatorname{im}A_v$ pointwise. If every one-site slice of
    $\psi$ at $v$ belongs to $\operatorname{im}A_v$, for all $v$,
    then $(\bigotimes_vP_v)\psi=\psi$.
    In particular, if $L_vF_v$ fixes $\operatorname{im}A_v$ and
    every vertex belongs to some $R_i$, then
    $L_{\mathrm{all}}F_{\mathrm{all}}\psi=\psi$ for every
    $\psi\in\mathcal P_{\mathcal R}(A)$.
\end{theorem}
\begin{proof}\leanok
    \uses{thm:peps_region_vertex_image_support}
    The operator acting as $P_v$ at one vertex and as the identity
    elsewhere fixes $\psi$, because it acts on its one-site slices.
    Multiplying these operators gives the first assertion.
    Vertex coverage puts every one-site slice of a parent vector in
    the appropriate site image. Take $P_v=L_vF_v$ and use
    $L_{\mathrm{all}}F_{\mathrm{all}}=\bigotimes_v(L_vF_v)$.
\end{proof}

\begin{theorem}[Equivalence of covered parent spaces]%
    \label{thm:peps_range_parent_equivalence}
    \lean{TNLean.PEPS.globalPhysicalMap_mem_parent_physicalDeform,
        TNLean.PEPS.map_regionParentGroundSpace_physicalDeform_of_leftInverse,
        TNLean.PEPS.map_regionParentGroundSpace_physicalDeform_of_injOn,
        TNLean.PEPS.globalPhysicalMap_injOn_parent,
        TNLean.PEPS.regionParentGroundSpaceRangeEquiv,
        TNLean.PEPS.finrank_regionParentGroundSpace_physicalDeform_of_injOn}
    \leanok
    Every physical product map satisfies
    $F_{\mathrm{all}}\mathcal P_{\mathcal R}(A)
        \subseteq\mathcal P_{\mathcal R}(B)$, without injectivity
    or coverage assumptions.
    If each $F_v$ is injective on $\operatorname{im}A_v$ and the
    regions cover all vertices, its restriction is a linear isomorphism
    \begin{align}
        F_{\mathrm{all}}:\mathcal P_{\mathcal R}(A)
                                       & \xrightarrow{\ \cong\ }\mathcal P_{\mathcal R}(B),
        \label{eq:peps_range_parent_equivalence}                                            \\
        \dim\mathcal P_{\mathcal R}(B) & =\dim\mathcal P_{\mathcal R}(A).
    \end{align}
    An inverse is $L_{\mathrm{all}}$ for any site inverses from
    Theorem~\ref{thm:peps_range_parent_site_inverse}.
\end{theorem}
\begin{proof}\leanok
    \uses{thm:peps_range_parent_site_inverse,
        thm:peps_range_parent_slice_reconstruction}
    The identity $\mathcal G_S(B)=F_S\mathcal G_S(A)$ sends each
    regional constraint to the corresponding constraint for $B$.
    In a fixed exterior slice, the physical maps on exterior vertices
    only form linear combinations of source slices; hence the forward
    inclusion holds for arbitrary $F_v$.
    Apply it also to $L_v$, using $LB=A$.
    Since $LF$ fixes the source site images and $FL$ fixes the target
    site images, Theorem~\ref{thm:peps_range_parent_slice_reconstruction}
    gives both inverse identities on the parent spaces.
    This proves~(\ref{eq:peps_range_parent_equivalence}) and the
    dimension equality.
\end{proof}

\begin{theorem}[Kernels of positive regional parents]%
    \label{thm:peps_range_parent_kernel}
    \lean{TNLean.PEPS.map_ker_regionParentHamiltonian_physicalDeform_of_injOn}
    \leanok
    \uses{def:peps_region_parent_interaction}
    Assume the site-image injectivity and vertex coverage of
    Theorem~\ref{thm:peps_range_parent_equivalence}, and suppose
    $\mathcal I$ is finite. Let $H_i,K_i$ be positive semidefinite
    regional interactions with
    $\ker H_i=\mathcal G_{R_i}(A)$ and
    $\ker K_i=\mathcal G_{R_i}(B)$.
    If $H$ and $K$ are the sums of their placements on the full graph,
    then $F_{\mathrm{all}}(\ker H)=\ker K$.
\end{theorem}
\begin{proof}\leanok
    \uses{thm:peps_range_parent_equivalence}
    Positivity identifies the kernel of each sum with the intersection
    of its local kernels. Thus $\ker H=\mathcal P_{\mathcal R}(A)$
    and $\ker K=\mathcal P_{\mathcal R}(B)$, to which
    (\ref{eq:peps_range_parent_equivalence}) applies.
\end{proof}

\begin{theorem}[Physical comparison for one fixed virtual representation]%
    \label{thm:peps_range_parent_g_comparison}
    \lean{TNLean.PEPS.IsGInjective.exists_physicalMap_comp_eq,
        TNLean.PEPS.exists_physicalDeform_eq_of_isGInjective}
    \leanok
    \uses{def:peps_g_injective}
    Let $G$ be finite, $\rho$ a complex representation on $\mathcal W$,
    and $T:\mathcal W\to\mathcal P$ and
    $S:\mathcal W\to\mathcal Q$ be $G$-injective for $\rho$.
    There is a linear $F:\mathcal P\to\mathcal Q$ such that
    $FT=S$ and $F$ is injective on $\operatorname{im}T$.
    In particular, two tensor families $A,B$ on the same virtual
    bond spaces, with site maps $G$-injective for the same $\rho_v$,
    admit such maps $F_v$ with $B_v=F_vA_v$.
    Their ambient physical dimensions may differ.
\end{theorem}
\begin{proof}\leanok
    \uses{thm:peps_g_injective_physical_map}
    Choose $L$ with $LT=\Pi_\rho$, where
    $\Pi_\rho=|G|^{-1}\sum_g\rho(g)$. Set $F=SL$.
    Invariance of $S$ gives $FT=S\Pi_\rho=S$.
    Since both $T$ and $FT=S$ are $G$-injective,
    Theorem~\ref{thm:peps_g_injective_physical_map} gives injectivity
    of $F$ on $\operatorname{im}T$.
    Apply this argument at every vertex and evaluate on virtual basis
    vectors to obtain equality of the tensor coefficients.
\end{proof}

\begin{theorem}[Parent kernels for a fixed G-injective representation]%
    \label{thm:peps_range_parent_g_transport}
    \lean{TNLean.PEPS.map_regionParentGroundSpace_physicalDeform_of_isGInjective,
        TNLean.PEPS.exists_regionParentGroundSpace_transport_of_isGInjective,
        TNLean.PEPS.exists_ker_regionParentHamiltonian_transport_of_isGInjective}
    \leanok
    Let $A,B$ have the same virtual bond spaces and let their site maps
    be $G$-injective for the same representations $\rho_v$ of a finite
    group. If the regions $R_i$ cover every vertex, there are rectangular
    physical maps $F_v$ such that, for every region $S$,
    \begin{align}
        F_S\mathcal G_S(A) & =\mathcal G_S(B),            \\
        F_{\mathrm{all}}\mathcal P_{\mathcal R}(A)
                           & =\mathcal P_{\mathcal R}(B).
    \end{align}
    The second restriction is injective, and the parent spaces have
    equal dimensions. If $B=FA$ is already specified and both site
    families are $G$-injective for $\rho_v$, the second equality holds
    for that $F$.
    For a finite region family and arbitrary positive regional
    interactions with these exact local kernels, the same comparison
    gives $F_{\mathrm{all}}(\ker H)=\ker K$, injectivity on
    $\ker H$, and $\dim\ker K=\dim\ker H$.
\end{theorem}
\begin{proof}\leanok
    \uses{thm:peps_range_parent_g_comparison,
        thm:peps_range_parent_equivalence,
        thm:peps_g_injective_physical_map}
    Theorem~\ref{thm:peps_range_parent_g_comparison} supplies
    $B=FA$ and site-image injectivity. Regional contraction commutes
    with $F$, giving the first equality. Apply
    Theorem~\ref{thm:peps_range_parent_equivalence} to the common
    parent spaces. For a prescribed $F$, its site-image injectivity
    follows directly from $G$-injectivity of $A$ and $FA$.
    Finally, positivity identifies each Hamiltonian kernel with its
    common parent space.
\end{proof}
% These fixed-representation consequences do not change virtual bond spaces.
% They are not a completion of SCP10 Theorem 5.7 for arbitrary semi-regular
% canonical parents, nor a regular/semi-regular parent-kernel comparison.
